% ============================================================
% 第 19 章  提高电力系统稳定性的措施
% （教材：何仰赞、温增银《电力系统分析》下册 第 19 章）
% 结构对齐 19.1—19.4；例题、习题不录
% ============================================================

\chapter{提高电力系统稳定性的措施}

前四章讲“怎样判断稳定”，本章讲“怎样提高稳定”。所有措施都可以从两个式子找到落点：
\[
  P_e=P_{\max}\sin\delta=\frac{EU}{X}\sin\delta
  \qquad\text{与}\qquad
  T_J\frac{\mathrm{d}^{2}\delta}{\mathrm{d}t^{2}}=P_T-P_e .
\]
前者说明“提高功率极限”的途径（提高 $E$、$U$，减小 $X$）；
后者说明“减小加速面积、增大减速面积、缩短故障持续时间”的途径。
本章按“元件—附加装置—运行方式”三个层次把这些途径展开。

\section{提高稳定性的一般原则}

由等面积定则（第 17 章），暂态稳定的条件是减速面积大于加速面积：
\begin{equation}
  A_{dec}=\int_{\delta_c}^{\delta_{\max}}\left(P_{\text{III}}-P_T\right)\mathrm{d}\delta
  >A_{acc}=\int_{\delta_0}^{\delta_c}\left(P_T-P_{\text{II}}\right)\mathrm{d}\delta .
\end{equation}
由此得到三条一般原则：
\begin{enumerate}
  \item \textbf{提高故障后的功率特性} $P_{\text{III}}$（增大减速面积）：
        提高电压、减小电抗、增加并联回路、投入强励；
  \item \textbf{减小故障期间的功率降落}（减小加速面积）：
        采用单相重合闸、快速切除、减小故障期间的等值电抗；
  \item \textbf{缩短故障持续时间}（减小加速面积）：
        加快继电保护动作、加快断路器开断、采用快速开关。
\end{enumerate}
此外还有两条辅助原则：
\begin{itemize}
  \item \textbf{增大惯性时间常数} $T_J$（减缓 $\delta$ 变化）；
  \item \textbf{增大阻尼}（抑制振荡，提高静态稳定与动态稳定）。
\end{itemize}

\begin{table}[htbp]
  \centering
  \caption{提高稳定性的三条基本途径}
  \small
  \begin{tabular}{@{}p{0.22\linewidth}p{0.30\linewidth}p{0.38\linewidth}@{}}
    \toprule
    途径 & 作用对象 & 典型措施 \\
    \midrule
    提高功率极限 $P_{\max}$ & $E$、$U$、$X$ & 强励、串补、分裂导线、双回线、中间开关站 \\
    减小加速面积 & 故障期间的 $P_{\text{II}}$、$\delta_c$ & 快速保护、快速开关、单相重合闸、电气制动 \\
    增大减速面积 & 故障后的 $P_{\text{III}}$ & 强励、快速励磁、并联补偿、切机切负荷 \\
    \bottomrule
  \end{tabular}
\end{table}

\subsection{三条途径的定量比较}

把等面积定则写成“临界切除角”的形式（第 17 章），可以定量看出各措施的相对效果：
\begin{equation}
  \cos\delta_{cm}
  =\frac{P_{\text{III}m}\cos\delta_{\max}-P_{\text{II}m}\cos\delta_0
  +P_T\left(\delta_0-\delta_{\max}\right)}{P_{\text{III}m}-P_{\text{II}m}} .
\end{equation}
\begin{itemize}
  \item \textbf{提高 $P_{\text{III}m}$}（强励、并联补偿、增加回路）：分母增大、分子中的
        $P_{\text{III}m}\cos\delta_{\max}$ 项也增大，$\cos\delta_{cm}$ 减小，
        $\delta_{cm}$ 增大——允许切除角变大，效果显著且“线性”；
  \item \textbf{提高 $P_{\text{II}m}$}（减小故障期间的电抗、单相重合闸）：分子中
        $-P_{\text{II}m}\cos\delta_0$ 项增大（$\cos\delta_0>0$ 时使分子减小）、
        分母增大，$\delta_{cm}$ 增大——效果也很显著，且不增加稳态设备投资；
  \item \textbf{缩短切除时间}：不改变 $\delta_{cm}$，而是使实际切除角 $\delta_c<\delta_{cm}$，
        相当于“把工作点移到安全区内”——这是最直接的措施，
        且保护与开关的投资远低于一次设备。
\end{itemize}
工程上的优先级通常是：\textbf{先加快切除（最便宜），再加强励与单相重合闸，
最后才考虑串补、并补等一次设备}。

\section{改善电力系统基本元件的特性和参数}

\subsection{改善发电机及其励磁调节系统的特性}

\begin{itemize}
  \item \textbf{快速励磁与强行励磁}：故障时迅速把励磁电压升到顶值（$1.8\sim2$ 倍额定），
        维持 $E_q$、提高 $P_{\text{III}}$，是最有效、最经济的措施之一；
  \item \textbf{装设电力系统稳定器（PSS）}：提供正阻尼，抑制低频振荡，
        使高放大倍数的快速励磁系统既能提高极限又不致振荡；
  \item \textbf{减小 $X_d$}：增大发电机短路比（减小同步电抗），可提高功率极限，
        但会增加电机尺寸与造价；
  \item \textbf{增大惯性时间常数 $T_J$}：转子加重，$\delta$ 变化变慢，有利于暂态稳定，
        但对调频不利。
\end{itemize}

\subsection{改善原动机的调节特性}

\begin{itemize}
  \item \textbf{快速关闭汽门（快关）}：故障时迅速减小汽轮机进汽量，降低 $P_T$，
        从而减小加速面积；故障切除后再快速恢复。这是火电机组提高暂态稳定的重要手段；
  \item \textbf{改善调速器特性}：使原动机出力能快速跟随，但调速器过快的响应可能
        引起机械功率振荡，须与系统阻尼协调；
  \item \textbf{水轮机的特点}：水锤效应使水轮机出力在导叶快速关闭时先增后减，
        调节不当会恶化稳定，故水轮机快关须谨慎。
\end{itemize}

\subsection{减小变压器的电抗}

增大变压器容量、降低短路电压百分数 $U_s\%$，可减小 $X_T$，
从而提高功率极限、减小电压损耗。但 $U_s\%$ 过小会使短路电流增大，
须与开关设备的开断能力协调。

\subsection{改善继电保护和开关设备的特性}

\begin{itemize}
  \item \textbf{快速切除故障}：缩短 $\delta_c$（极限切除角不变，切除越早越安全），
        是最直接有效的措施；
  \item \textbf{单相自动重合闸}：单相接地故障时只切除故障相，其余两相继续运行，
        保持了线路的功率传输能力，显著减小加速面积；
  \item \textbf{综合重合闸}：根据故障类型自动选择单相或三相重合，兼顾可靠性与稳定性；
  \item \textbf{快速开关设备}：缩短开断时间（从 $3\sim5$ 周波降到 $2$ 周波）。
\end{itemize}

\subsection{改善输电线路的特性}

\begin{itemize}
  \item \textbf{采用分裂导线}：减小线路电抗（$x_1$ 从 $0.4$ 降到 $0.28\sim0.33\ \Omega$/km），
        提高功率极限；
  \item \textbf{串联电容补偿}：直接抵消部分线路电抗，是提高远距离输电能力最有效的手段；
  \item \textbf{增大导线截面}：减小电阻与电抗（电抗变化不大），主要降低网损；
  \item \textbf{增加并联回路数}：双回线使等值电抗减半，可靠性也提高。
\end{itemize}

\subsection{采用直流输电}

直流输电在提高稳定性方面有独特优势：
\begin{itemize}
  \item 直流联络线\textbf{不传输同步功率}，两端交流系统不必保持同步，
        从根本上消除了“联络线失步”问题；
  \item 直流输电可以\textbf{快速调节功率}（毫秒级），故障时迅速降低输送功率，
        减轻送端加速、受端功率缺额；
  \item 直流输电可以\textbf{隔断}两侧系统的故障传递，避免连锁反应；
  \item 直流线路本身无电抗压降、无充电功率，适合远距离大容量输电。
\end{itemize}

\section{采用附加装置提高电力系统的稳定性}

\subsection{输电线路采用串联电容补偿}

在线路中串入电容器，等值电抗变为 $X-X_C$，功率极限提高为
\begin{equation}
  P_{\max}=\frac{EU}{X-X_C}=\frac{EU}{X\left(1-K_C\right)},
  \qquad K_C=\frac{X_C}{X}.
\end{equation}
补偿度 $K_C$ 一般取 $0.3\sim0.5$（过高会引起次同步谐振与过电压）。
串补的附加好处是减小电压损耗、提高输电能力；须配置的保护有：
电容器过电压保护、间隙保护、旁路开关。

\subsection{输电线路的并联电抗补偿}

并联电抗器主要用于\textbf{限制空载或轻载时的工频过电压}（长线电容效应），
它吸收容性无功、降低电压。虽然并联电抗本身降低电压、不利于功率传输，
但在长距离超高压线路中是必需的；运行时可按需投切。
与并联电抗相反，\textbf{并联电容器与静止补偿器}可提高电压、增大功率极限，
是提高稳定性的有效手段（尤其在受端系统）。

\subsection{输电线路设立开关站}

长距离线路中途设立开关站，把线路分成若干段：
\begin{itemize}
  \item 故障时只切除一段，其余段仍可运行，减小了故障影响范围与加速面积；
  \item 开关站可装设并联补偿装置，改善电压分布；
  \item 相当于缩短了电气距离，提高了功率极限。
\end{itemize}

\subsection{中继同步调相机}

在线路中途装设同步调相机，可维持中间点电压、提供无功支撑，
从而提高功率极限（相当于把长线分成两段、中间电压恒定）。
现代系统中多由静止补偿器（SVC／STATCOM）取代，因后者响应更快、维护更简单。

\subsection{变压器中性点经小阻抗接地}

单相接地短路时，故障电流经中性点接地阻抗流通。在变压器中性点串入小阻抗 $Z_n$，
相当于在零序网中串入 $3Z_n$，增大了零序阻抗、减小了单相接地短路电流，
从而减小故障期间的功率降落（减小加速面积）。这一措施对以单相接地为主的
超高压系统效果明显，但须注意它对零序保护整定的影响。

\subsection{发电机采用电气制动}

故障发生时，迅速投入接在发电机端的\textbf{制动电阻}，把转子的加速动能转化为热能消耗掉，
相当于人为增大减速面积。制动电阻在故障切除后退出。
电气制动的优点是响应快（毫秒级）、可控；缺点是需附加设备与散热措施，
且制动过强可能引起过制动（反向失步）。

\begin{table}[htbp]
  \centering
  \caption{附加装置类措施的比较}
  \small
  \begin{tabular}{@{}p{0.20\linewidth}p{0.24\linewidth}p{0.22\linewidth}p{0.26\linewidth}@{}}
    \toprule
    措施 & 主要作用 & 响应速度 & 代价／限制 \\
    \midrule
    串联电容补偿 & 减小电抗、提高 $P_{\max}$ & 固定（不可调） & 次同步谐振、过电压保护 \\
    并联电抗器 & 限制工频过电压 & 可投切 & 降低电压，不利于输送 \\
    并联电容器／SVC & 提高电压与 $P_{\max}$ & 快（SVC 毫秒级） & 投资、谐波 \\
    中间开关站 & 缩小故障范围 & — & 占地、投资 \\
    中性点小阻抗 & 减小单相短路电流 & — & 影响零序保护 \\
    电气制动 & 消耗加速动能 & 毫秒级 & 设备与散热，可能过制动 \\
    \bottomrule
  \end{tabular}
\end{table}

\section{改善运行条件及其他措施}

\subsection{正确规定电力系统运行参数的数值}

\begin{itemize}
  \item \textbf{运行角}：正常运行时功角应留有足够裕度（$K_p\ge15\%\sim20\%$），
        避免长期在接近 $90^\circ$ 处运行；
  \item \textbf{电压水平}：维持较高的运行电压（但不超过设备允许值），
        可提高功率极限；
  \item \textbf{备用容量}：留足旋转备用，事故时可快速支援；
  \item \textbf{线路潮流}：控制重载线路的输送功率，避免任一元件的故障引起连锁过载。
\end{itemize}

\subsection{合理选择电力系统的运行接线}

\begin{itemize}
  \item 避免\textbf{弱联}（单回长线联络），必要时增设第二回线或改用直流联络；
  \item 合理配置\textbf{环网}，使任一元件退出后潮流仍能合理分布；
  \item 重要负荷采用双电源供电，重要断面避免“卡脖子”；
  \item 检修方式下须重新校核稳定，必要时限制输送功率。
\end{itemize}

\subsection{切除部分发电机及部分负荷}

\begin{itemize}
  \item \textbf{送端切机}：故障后送端功率过剩、机组加速，快速切除部分机组可减小加速面积，
        是最有效的紧急控制措施之一；
  \item \textbf{受端切负荷}：受端功率缺额、频率下降时，切除部分负荷可维持系统平衡，
        防止频率崩溃；
  \item 切机切负荷须与保护、安控装置配合，动作要快（$100\sim200$ ms 级）。
\end{itemize}

\subsection{高压直流输电功率的快速调节}

直流输电的功率可以快速、独立地调节（毫秒级）。故障时：
\begin{itemize}
  \item 送端直流功率快速降低，减小送端加速；
  \item 受端直流功率快速提升（若有备用），支援受端；
  \item 直流功率调制还可为交流系统提供附加阻尼，抑制低频振荡。
\end{itemize}

\subsection{减少系统稳定破坏所带来的损失和影响}

稳定破坏（失步、振荡、电压崩溃）一旦发生，须尽量缩小影响范围：
\begin{enumerate}
  \item \textbf{解列}：在预定地点（解列点）把系统分成若干独立部分，
        防止事故扩大；解列点应使各部分功率尽量平衡；
  \item \textbf{低频减载}：频率下降到整定值时按轮次切除负荷，防止频率崩溃；
  \item \textbf{振荡解列}：检测到失步振荡时在振荡中心解列；
  \item \textbf{黑启动}：事故后依靠具备自启动能力的机组恢复供电。
\end{enumerate}

\begin{table}[htbp]
  \centering
  \caption{提高稳定性措施的分类总结}
  \small
  \begin{tabular}{@{}p{0.14\linewidth}p{0.36\linewidth}p{0.40\linewidth}@{}}
    \toprule
    层次 & 措施 & 特点 \\
    \midrule
    元件参数 & 减小 $X$（分裂导线、双回线、低 $U_s\%$ 变压器）、增大 $T_J$、
      快速励磁与强励、PSS & 一次投资，长期有效 \\
    附加装置 & 串补、并补、开关站、调相机、中性点小阻抗、电气制动 & 针对性强，需附加设备 \\
    运行方式 & 控制运行角与潮流、合理接线、留备用 & 不花钱，效果有限 \\
    紧急控制 & 快速保护、单相重合闸、切机切负荷、直流调制、解列 & 事故时动作，是“最后一道防线” \\
    \bottomrule
  \end{tabular}
\end{table}

\corepoint{本章是全书最“工程”的一章，答题时不必罗列全部措施，而要\textbf{按机理归类}：
凡提高 $P_{\max}$ 的（减小 $X$、提高 $E$ 与 $U$）归一类；
凡减小加速面积的（快速切除、单相重合闸、快关汽门、电气制动）归一类；
凡增大减速面积的（强励、并补）归一类；
凡防止事故扩大的（解列、低频减载）归一类。
再补一句“各措施的经济性与适用场合”，答案就完整了。}
